\honest{Uncritical Supercriticality}{Eliezer Yudkowsky}{2007}

People argue over whether atheism is a ``religion.'' Arguing over the meaning of a word ``nearly always means that you've lost track of the original question.'' An atheist blames religion for the Inquisition and the Crusades. The believer answers that atheism is a religion too (``If atheism is a religion, then not collecting stamps is a hobby''), or that Stalin was an atheist (``Stalin's religion was Communism''; ``If Communism is a religion, then Star Wars fandom is a government''). I show with two made-up definitions that whether Stalin counts as religious depends only on the definition. He is religious if that means a definite opinion about some god, below 10\% or above 90\%. He is not if it means a probability above 90\% that a god exists. Redefining a word does not change history. The real question is why large groups of people have been slaughtered in the name of an idea.

Communism was a complex catastrophe, and there may be no single cause. But ``if I had to suggest an ur-mistake,'' God states it in Deuteronomy 13: if your brother, your child, your spouse or your closest friend tempts you to other gods, ``you must not listen to him'' and ``you must kill him.'' Stalin and Hitler set the same rule: do not argue with the critic, turn him in to the secret police. \nb{The record bears this out in substance. The Soviet state made a hero of Pavlik Morozov, a boy said to have denounced his father, and most Gestapo investigations began with denunciations by ordinary Germans. That shows the rule was present where the killing happened, not that it came first and caused the killing.}

Remember the remedy from my last post: question each added nice claim, so that one claim triggers fewer than one more. The opposite is when it feels wrong to argue against any positive claim at all. Arguments are soldiers. If you feel that correcting a flawed argument for evolution would help the creationists, or that you earn spiritual credit for each nice thing you say about God, or that the room will turn on you for ``not supporting our troops,'' or that saying anything against Communism gets you ``stoned to death shot,'' the spiral has gone supercritical. \nb{In 2007 the words ``stoned to death'' were struck through. The strikethrough has been lost in the text LessWrong now shows.}

So religion is not the key category. The best distinction I have heard between supernatural and naturalistic views is that the supernatural asserts ``ontologically basic mental substances.'' \nb{The distinction is Richard Carrier's, from January 2007. Yudkowsky named him in a later post, ``Excluding the Supernatural.''} Supernatural claims ``always turn out to be wrong'' for fairly fundamental reasons, but blaming them alone buys into religious exceptionalism. Spirals form easily around monotheisms whose God is defined by agreeing with every nice statement, especially once they threaten punishment for disbelief. But they also form around political systems, leaders, racial destiny and economic theories. Sam Harris came closer by blaming faith. New Agers inherit from Christianity the idea that faith is good, but not the exclusive scripture that keeps out rival ideas, so they spiral around ``stars, trees, magnets, diets, spells, unicorns.'' I give no evidence about Christianity or the New Agers.

The spiral ``turns much deadlier after criticism becomes a sin, or a gaffe, or a crime.'' \nb{Here the suggestion of the ur-mistake paragraph has become a plain statement, with no new evidence.} Some things deserve great praise, so you cannot flatly forbid praise beyond some point. But there is ``never an Idea so true that it's wrong to criticize any argument that supports it.'' Most possible beliefs are false, and so are most possible arguments for a true belief. And it is ``triple ultra forbidden'' to answer criticism with violence. ``Bad argument gets counterargument. Does not get bullet.''
